% Nouveau EGA V (cote 25): archive cote 25, site page 57 (PDF page 58).
% typescript. Draft compared: none.
% AI-transcribed by Claude Opus 5.5 (claude-opus-5-5, adaptive thinking, effort high), Message Batches API, batch ega25b (msgbatch_01A7n5YuXubUPNK1zsepi8RG), 2026-11-13; not yet checked by eye.
\scanpage{25/57}
\mspage{\unsure{92}}

L'hypothèse  dim $f(X)\geqslant 2$  signifie que  le degré de transcendance de
$k(t_1,..,t_r)$ sur $k$  est $\geqslant 2$, i.e. on peut trouver  $t_1,t_2$ disons alg ind.
Dans a) on prendra alors   $P = T_1$ , $Q = T_1^{\,2}$ , $R = T_1T_2$ , Dans b), on
fait pareil, \xout{\unsure{mais en ayant soin de prendre} \unsure{$t_1$}} en notant qu'on peut ci-%
dessus choisir $t_1$ provenant bien de $T_1$ \uline{nul sur} L, grâce au fait $f(X)\not\subset L$
(qui implique qu'il  existe un indice  i  entre 1 et s tel que $t_s$ non
nul, donc $t_s$ non constante (car $t_s$ s'annule en x) donc $t_s$  non alg sur k
( NB on peut supposer k alg clos ).  Le cas c) résulte de b), \deleted{avec n=2
même au lieu de n=3,} sauf dans le cas où \xout{X} f(X) est contenu dans
l'image \inserted{L} par f de l'espace tangent à X en x, \deleted{ce qui  exige d'ailleurs pou
des raisons de dim que l'application tangente en x est injective i.e.
f non ramifié en x, \xout{Lorsque dim X} de sorte que f(X) est contenu dans
une sous-variété linéaire de P de même dimension que X.} \hand{[}Lorsque dim X=2
ce cas est effectivement exceptionnel (la trace d'une quadrique tangente
à un plan sur ce plan est en général formé de deux droites concourran-%
tes, donc n'est pas irréductible)\hand{]}. Mais pour traiter ce cas, \xout{sous} la
forme P,Q,R explicitée plus haut, on peut remplacer \inserted{si dim L $=2$} évidamment X par
L lui-même, où la solution est triviale ( prendre $P=T_r^{\,2}$\xout{\illegible{1}}, $Q=T_r^{\,3}$ ,
$R = T_r^{\,2}T_{r-1}$ , en notant que  $Q/P = T_r$ et  $R/P = T_{r-1}$ sont des formes
linéaires lin ind sur L, donc alg ind\inserted{si dim $L\geqslant 3$, alors $T_{r-2}$, $T_{r-1}$, $T_r$ lin. indép. sur L, et on prendra $P=T_r^{\,2}$, $Q=\unsure{T_{r-1}^{\,2}}$, $R=\unsure{T_{r-2}^{\,2}}$}). \deleted{On pourrait signaler en
remarque que dans c), on peut même prendre n=2  sauf si  f(X) est conte-%
nu dans un plan .}

Conjuguant avec \unsure{8.181|8.18} on trouve un

\uline{Corollaire} 9.4.   \xout{(Pour X propre sur k}

Enfin, il faut encore conjuguer  ce dernier avec 9.2. pour
trouver un théorème récapitulatif  du ``cas excellent'':
\deleted{\uline{Théorème}} 9.5.   (Lorsque X est lisse \deleted{et} propre\inserted{(et géom intègre)} sur k, et $f:X\to P$ non
ramifié,\hand{\unsure{X}}de dimension $\geqslant 2$)\hand{, en signalant le résultat sur \unsure{$\pi_1$} ?.}
\hand{quand $X\to P$ non immersion (variété des pts singuliers de \unsure{$f$}).}

\margin{57}
